% !TeX root = ../papers/nc-p3.tex
\section{五つの水平台の中心束と第一長方形の階数十四曲線の全正則性}
\label{sec:rectangle-quartic-bundles}
\begin{theorem}\label{thm:rectangle-quartic-bundles}
第二長方形の水平台143から147で $b+c\ne0$ なら、四次中心は四元群の文字順に
階数 $(3,2,0,0)$ のベクトル束をなし、各ファイバーでの中心と可換する。
第一長方形の台85の反対階数14曲線では、全点で四中心四次式上自由階数256の
標準Koszul左右AS正則代数が得られる。全点が最小周期4で、指定始点の無印同型類は
定理\ref{thm:first-rectangle-rank-strata}の $r$ で一意に決まる。
\end{theorem}
\begin{proof}[第二長方形の中心を解く共通表示]
各位置の文字番号に順に $(0,0,2,2)$ を加える。全和が零なので四次の文字は変わらない。
$b\ne1$ の表示 $t,b,x=a,y$ で $c=(t^2-b)/(1-b)$ とし、左上行列式の単元因子を
$E=(b+c)xy-bc(x+y)+b^2(c-1)$ と書く。
$tbc(b-1)(b+c)E$ および $q=x,y$ に対する $q(q-b)(q-1)(cq+b(1-c))$ の分子・分母の因子を逆にする。
これは定理\ref{thm:horizontal-rectangle-low-fibers}の共通の開条件で、五つの台で消えてよい座標を要求しない。
各二次関係のフレームは可逆小行列式を持つ。各始点の長さ5までの経路と、四色の長さ4の経路を一文字ずつ左右から
延長して差を取る可換方程式を、この基底環の単元だけで掃き出せる。経路の次元列は
$1,4,10,20,35,56$、可換方程式の自由列の数は $3,2,0,0$ であり、核・余核は特殊化と可換である。

$b=1,c\ne1$ は一位置移動でこの表示に入る。文字の線型部は厚い二色の差である文字1を保ち、
位置別平行移動の総和は零なのでこの移動でも四次の文字0,1を保つ。
$b=c=1$ では $t=1,-1$ を個別に代入する。今度は $E=2xy-x-y$ と
$x(x-1)y(y-1)$ を逆にした二変数環で、同じ平行移動が両符号で再び単元だけの掃き出しを与える。
\path{computations/verify_second_rectangle_quartic_bundle.py} が三表示を検算する。
\end{proof}

\begin{proof}[第一長方形の曲線を全点で覆う]
$u=r(r-1)/(r+1)$, $v=(1-r)/(r(r+1))$, $z=-4v/(v+1)^2$ を使う。
始点0の五中心式上の表示に定理\ref{thm:first-rectangle-edge-regularity}の独立関係による判定を適用する。
今回は全中心商の基底個数は長さ0から8まで
$(1,4,10,20,30,36,34,20,0)$ となる。従って必要な新関係は次数9から16まで順に
$20,34,36,30,20,10,4,1$ 個である。中心変数の定数評価と位置別文字移動を
$((1,0,0,0,0),(0,2,2,0))$ および $((0,1,0,0,0),(2,0,2,0))$ と取る。
元の曲線で可逆な因子を除けば、中心核・商基底・全ての関係の独立性を同時に保証する二表示の
例外既約因子は120個と141個である。その共通部分は次の八数体で尽くされる。
\[
\begin{gathered}
r^2+6r+1,\quad r^2-3,\quad r^2-6r+1,\quad r^2-6r-1,\\
r^4\pm4r^3+2r^2\mp12r+9,\qquad
r^4\pm4r^3-6r^2\mp12r+1 .
\end{gathered}
\]
それぞれで定義域を保つ次数1の単純剰余根を取り、四中心商を長さ13まで全始点で検算する。
\path{research/first_rectangle_rank_fourteen_residues.json} に素数と根を、
\path{computations/verify_first_rectangle_rank_fourteen_curve.py} に二表示・八数体の再現計算を掲げる。
曲線は反対階数で内在的に定まるため一位置移動で保たれる。従って全点で始点0の表示が得られれば、
全点の全始点で自由性と深さ4が従う。四始点に共通の一般の四中心式を取り左右正則性を得る。
定理\ref{thm:first-rectangle-horizontal-frames}の周期2の二直線とは交わらない。
\end{proof}
